\chapter{完整源码清单：Compute Systems 并发实验}

本附录完整收录 \filepath{books/cpu-volume-2/labs/compute_systems}。第二册主线工程先用 chapter models 给每章建立可运行模型，再从串行 reference、并行归约、同步原语、有界队列、wait channel、任务运行时一路推进到 futex 合同。阅读时不要跳过测试：系统代码最重要的证据往往在失败路径测试里。

\section{一、CMakeLists.txt}

CMake 文件把库、命令行 probes 和测试目标接在一起。读者应看清每个源文件进入哪个 target，以及 Linux 平台相关代码如何被组织。

\lstinputlisting[language=bash,caption={compute\_systems/CMakeLists.txt}]{../../labs/compute_systems/CMakeLists.txt}

\section{二、README}

README 给出实验范围和运行命令。它说明本目录是第二册原理卷的本地实验层，完整工程化恢复和 manifest 逻辑会继续衔接代码实践卷。

\lstinputlisting[language=bash,caption={compute\_systems/README.md}]{../../labs/compute_systems/README.md}

\section{三、逐章模型头文件}

\filepath{chapter_models.hpp} 为第二册每个主章和深入章定义一个报告结构和一个模型函数。它的目标是让“分支预测、TLB、NUMA、Roofline、背压、epoch、quorum、checkpoint”等词都有能编译的输入和输出。

\lstinputlisting[language=C++,caption={include/compsys/chapter\_models.hpp}]{../../labs/compute_systems/include/compsys/chapter_models.hpp}

\section{四、并行归约头文件}

\filepath{parallel_reduce.hpp} 定义串行 reference、并行归约和报告结构。读者应从这里确认每个元素必须恰好参与一次归约，worker 数只是调度参数，不改变语义。

\lstinputlisting[language=C++,caption={include/compsys/parallel\_reduce.hpp}]{../../labs/compute_systems/include/compsys/parallel_reduce.hpp}

\section{五、同步原语头文件}

这个头文件暴露 mutex counter、atomic counter 和 bounded queue。它把“线程安全”拆成可测试接口：add、value、push、pop、close 和统计报告。

\lstinputlisting[language=C++,caption={include/compsys/sync\_primitives.hpp}]{../../labs/compute_systems/include/compsys/sync_primitives.hpp}

\section{六、等待通道头文件}

wait channel 头文件描述 eventfd/epoll 合同报告。读者要区分“被唤醒”和“数据已经 drain 完成”，这两个事实不能混写。

\lstinputlisting[language=C++,caption={include/compsys/wait\_channels.hpp}]{../../labs/compute_systems/include/compsys/wait_channels.hpp}

\section{七、任务运行时头文件}

\filepath{task_runtime.hpp} 是第二册主线工程的运行时接口。它把提交、完成、失败、拒绝、队列深度和 CSV 报告都变成结构化字段。

\lstinputlisting[language=C++,caption={include/compsys/task\_runtime.hpp}]{../../labs/compute_systems/include/compsys/task_runtime.hpp}

\section{八、Futex 头文件}

futex 头文件只暴露合同 probe 和报告字段。读者应注意：futex 不是普通 mutex 的替代模板，它要分别证明快路径、慢路径和错误返回。

\lstinputlisting[language=C++,caption={include/compsys/futex\_lab.hpp}]{../../labs/compute_systems/include/compsys/futex_lab.hpp}

\section{九、逐章模型实现}

\filepath{chapter_models.cpp} 是对“每章完整模型代码”的直接落实。每个函数都从一个小输入出发，计算报告字段；测试文件会逐章断言这些字段，probe 会把报告打印出来。

\lstinputlisting[language=C++,caption={src/chapter\_models.cpp}]{../../labs/compute_systems/src/chapter_models.cpp}

\section{十、并行归约实现}

实现文件展示从串行 reference 到多 worker 切片的完整过程。阅读重点是切片边界、空输入、worker 数过大、局部和汇总和报告字段。

\lstinputlisting[language=C++,caption={src/parallel\_reduce.cpp}]{../../labs/compute_systems/src/parallel_reduce.cpp}

\section{十一、同步原语实现}

同步实现覆盖 mutex、atomic 和有界队列。读者应重点看 close 如何唤醒等待中的生产者和消费者，以及报告字段如何记录等待、关闭和队列水位。

\lstinputlisting[language=C++,caption={src/sync\_primitives.cpp}]{../../labs/compute_systems/src/sync_primitives.cpp}

\section{十二、等待通道实现}

wait channel 实现把 eventfd、epoll、队列和信号量语义放在同一个 probe 里。它的价值是把内核可读事件和用户态数据结构 drain 分开验证。

\lstinputlisting[language=C++,caption={src/wait\_channels.cpp}]{../../labs/compute_systems/src/wait_channels.cpp}

\section{十三、任务运行时实现}

任务运行时实现负责 worker 生命周期、任务队列、异常捕获、shutdown、同池等待 probe、continuation probe 和 CSV 导出。读者应逐段检查每个报告字段在什么路径上变化。

\lstinputlisting[language=C++,caption={src/task\_runtime.cpp}]{../../labs/compute_systems/src/task_runtime.cpp}

\section{十四、Futex 实现}

futex 实现覆盖用户态快路径、内核 wait/wake、EAGAIN、timeout 和 EINTR。读者应把这些分支和正文中“快路径不是全部合同”的说明对照。

\lstinputlisting[language=C++,caption={src/futex\_lab.cpp}]{../../labs/compute_systems/src/futex_lab.cpp}

\section{十五、主程序}

\filepath{src/main.cpp} 是最小命令行入口，用来确认库可以被普通可执行程序调用。它不承担完整测试责任，也不应该复制库内部逻辑。

\lstinputlisting[language=C++,caption={src/main.cpp}]{../../labs/compute_systems/src/main.cpp}

\section{十六、逐章模型 probe}

\filepath{chapter_models_probe.cpp} 运行 21 个逐章模型并打印稳定摘要。它让读者不需要先读完整测试，就能看到每章模型到底输出哪些字段。

\lstinputlisting[language=C++,caption={src/chapter\_models\_probe.cpp}]{../../labs/compute_systems/src/chapter_models_probe.cpp}

\section{十七、wait channel probe}

\filepath{wait_channels_probe.cpp} 把 eventfd/epoll 合同报告打印出来。读者应重点看输出字段是否能区分 ready、empty、FIFO 和 semaphore drain，而不是只看程序是否退出。

\lstinputlisting[language=C++,caption={src/wait\_channels\_probe.cpp}]{../../labs/compute_systems/src/wait_channels_probe.cpp}

\section{十八、TaskRuntime probe}

\filepath{task_runtime_probe.cpp} 运行任务运行时样本并输出报告。它用于人工观察 submitted、completed、failed、rejected 和 queue depth，不替代 \filepath{compsys_tests.cpp} 的断言。

\lstinputlisting[language=C++,caption={src/task\_runtime\_probe.cpp}]{../../labs/compute_systems/src/task_runtime_probe.cpp}

\section{十九、futex probe}

\filepath{futex_lab_probe.cpp} 打印 futex 合同样本。读者要把快路径计数、慢路径等待、唤醒和错误返回分开看，不能只用最终计数判断正确。

\lstinputlisting[language=C++,caption={src/futex\_lab\_probe.cpp}]{../../labs/compute_systems/src/futex_lab_probe.cpp}

\section{二十、完整测试}

测试文件是实验的核心证据。它不仅检查最终计数，还逐章断言 chapter models，并检查 close 唤醒、队列 drain、同池等待、CSV 表头、futex 错误路径等容易被漏掉的边界。

\lstinputlisting[language=C++,caption={tests/compsys\_tests.cpp}]{../../labs/compute_systems/tests/compsys_tests.cpp}
